\documentclass[11pt,a4paper]{article}
\usepackage[T1]{fontenc}
\usepackage{lmodern}
\usepackage[margin=26mm]{geometry}
\usepackage{amsmath,amssymb,amsthm,mathtools,mathrsfs}
\usepackage{microtype,booktabs,longtable,array}
\usepackage{xcolor,enumitem,fancyhdr,needspace}
\usepackage{xurl,hyperref,bookmark}
\definecolor{Ink}{HTML}{233240}
\definecolor{Accent}{HTML}{225C70}
\hypersetup{colorlinks=true,linkcolor=Accent,citecolor=Accent,urlcolor=Accent,
 pdftitle={Derivations, Differential Equations, and the Phase Obstruction over Surreal and Surcomplex Numbers},
 pdfauthor={Research exposition prepared for the Surreal documentation project}}
\setlength{\headheight}{14pt}
\setlength{\emergencystretch}{2em}
\pagestyle{fancy}\fancyhf{}
\fancyhead[L]{\small Surreal differential algebra}
\fancyhead[R]{\small Surreal / surcomplex numbers}
\fancyfoot[C]{\thepage}
\renewcommand{\sectionmark}[1]{\markboth{\thesection.\ #1}{}}
\numberwithin{equation}{section}
\newtheorem{theorem}{Theorem}[section]
\newtheorem{proposition}[theorem]{Proposition}
\newtheorem{lemma}[theorem]{Lemma}
\newtheorem{corollary}[theorem]{Corollary}
\theoremstyle{definition}
\newtheorem{definition}[theorem]{Definition}
\newtheorem{example}[theorem]{Example}
\theoremstyle{remark}
\newtheorem{remark}[theorem]{Remark}
\newcommand{\No}{\mathbf{No}}
\newcommand{\K}{\mathbb K}
\newcommand{\C}{\mathbb C}
\newcommand{\R}{\mathbb R}
\newcommand{\Q}{\mathbb Q}
\newcommand{\Z}{\mathbb Z}
\newcommand{\N}{\mathbb N}
\newcommand{\OO}{\mathcal O}
\newcommand{\mm}{\mathfrak m}
\newcommand{\PP}{\mathbf P}
\newcommand{\HH}{\mathscr H}
\newcommand{\EE}{\mathcal E}
\newcommand{\LL}{\mathcal L}
\newcommand{\DD}{\widetilde D}
\newcommand{\ct}{\operatorname{ct}}
\newcommand{\st}{\operatorname{st}}
\newcommand{\supp}{\operatorname{supp}}
\newcommand{\Ree}{\operatorname{Re}}
\newcommand{\Imm}{\operatorname{Im}}
\newcommand{\im}{\operatorname{im}}
\newcommand{\fin}{\operatorname{fin}}
\newcommand{\pipart}{\operatorname{pi}}
\newcommand{\Span}{\operatorname{span}}
\newcommand{\abs}[1]{\lvert #1\rvert}
\newcommand{\repo}{\texttt{VladimirReshetnikov/Surreal}}
\newcommand{\arxiv}[1]{\href{https://arxiv.org/abs/#1}{arXiv:#1}}
\newcommand{\code}[1]{\texttt{#1}}
\newcommand{\Dlog}{\ell_D}
\newcommand{\PS}{\mathbb S}
\newcommand{\Torus}{\mathbb T}
\newcommand{\val}{v}
\setlist[enumerate]{itemsep=3pt,topsep=5pt}
\setlist[itemize]{itemsep=3pt,topsep=5pt}

\begin{document}
\begin{titlepage}
\centering
\vspace*{12mm}
{\small\sffamily\color{Accent} A COMPANION TO THE SURREAL DOCUMENTATION PROJECT\par}
\vspace{12mm}
{\LARGE\bfseries\color{Ink} Derivations, Differential Equations,\par}
\vspace{3mm}
{\LARGE\bfseries\color{Ink} and the Phase Obstruction\par}
\vspace{5mm}
{\Large over Surreal and Surcomplex Numbers\par}
\vspace{9mm}
{\large From the Berarducci--Mantova derivation\par
 to Hahn workspaces and analytic differentiation\par}
\vspace{13mm}
\begin{minipage}{0.91\textwidth}
\small
\textbf{Abstract.}
A derivative of a function on the surreal numbers and a derivation of the
surreal scalar field are different objects. We develop the intrinsic
calculus associated with the normalized Berarducci--Mantova derivation
$D\omega=1$ and explain its relation to the analytic layers already documented
in the Surreal repository. Although $D$ is surjective on
$\K=\No[i]$, nonzero solutions of $Dy=qy$ need not exist. We prove an exact
criterion: for $q=a+ib$, such a solution exists precisely when a real
primitive of $b$ is finite. The obstruction is the purely infinite part of
that primitive. This gives a concrete exact sequence, a maximal domain for
an intrinsically compatible complex exponential, and an explicit
classification of constant-coefficient homogeneous equations. We also compute
both derivative images and logarithmic-derivative images in
$\C((\omega^{-\Q}))$, construct derivation-stable set-sized Hahn workspaces,
prove a two-derivative evaluation rule, and exhibit a Picard--Vessiot extension
that cannot embed differentially into $\No[i]$. Deep structural inputs are
identified; the deductions used to connect them are proved in detail.
\end{minipage}
\vfill
{\small Research exposition prepared with AI assistance\par}
\vspace{3mm}
{\small 21 September 2026\par}
\vspace{5mm}
{\footnotesize Repository snapshot inspected:\par}
{\footnotesize\ttfamily e260237db9b71da8b74a0c13c8e6355119091100\par}
\vspace{2mm}
{\footnotesize This article is not a refereed publication or a machine-checked proof.\par}
\end{titlepage}
\pagenumbering{roman}
\setcounter{tocdepth}{2}
\tableofcontents
\clearpage
\pagenumbering{arabic}

\section{The documentation gap and the question it raises}
\label{sec:gap}

The repository already contains substantial treatments of surreal algebra,
Hahn summation, surcomplex analytic geometry, contour functionals,
trigonometry, foundations, and exact computation. Another survey of those
subjects would add little. The gap addressed here is more specific:
\emph{what changes when the surreal coefficient field itself is equipped with
an intrinsic derivation?}

The inspected catalogue lists fifteen research packages. The foundations
README explicitly calls for separate interfaces for fine derivatives,
formal-series germs, coherent Hahn data, and scalar-field derivations.
The trigonometry README explicitly leaves a derivation on $\No$ unchosen
and distinguishes its phase-extension classification from a classification
of differential-equation solutions. These are sound boundaries, not defects.
What is missing from that documented interface is a worked theory linking the
two sides, with examples that prevent the boundaries from being erased
accidentally \cite{RepoMap,RepoFound,RepoTrig}.

\subsection{Scope of the repository audit}

The gap assessment is based on the catalogue, the READMEs of the analysis,
trigonometry, foundations and computer-algebra packages, and the opening
foundational section of the analysis source, all at the pinned revision on
the title page. It is not an exhaustive full-text audit of every source
manuscript, historical revision, or branch. In particular, this article does
\emph{not} claim that the repository never mentions Berarducci--Mantova,
transseries, or derivations. It supplies a dedicated development of a bridge
that the inspected documents distinguish but do not themselves claim to
complete \cite{RepoAnalysis,RepoCAS}.

\begin{center}
\small
\begin{tabular}{@{}p{0.22\textwidth}p{0.33\textwidth}p{0.36\textwidth}@{}}
\toprule
Existing package & Documented concern & Addition in this article\\
\midrule
Analysis & Functions, their germs, and Hahn-supported coefficients & Which derivative acts on an argument, and which acts on a scalar?\\[3pt]
Trigonometry & Finite phases and global phase conventions & Which phases can satisfy an intrinsic differential law?\\[3pt]
Foundations & Size and distinct mathematical interfaces & A support-level construction of $D$-stable Hahn workspaces\\[3pt]
Computer algebra & Exact representation and capability contracts & Exact differential tests in a rational Hahn workspace\\
\bottomrule
\end{tabular}
\end{center}

\subsection{The central contrast}

Fix the Berarducci--Mantova derivation $D$ with $D\omega=1$. Its extension to
$\K=\No[i]$ has constant field $\C$ and is surjective. Nevertheless,
\begin{equation}\label{eq:first-surprise}
 Dy=iy \quad\Longrightarrow\quad y=0.
\end{equation}
This is not a statement about failure to find a convenient notation for
$e^{i\omega}$. It is a nonexistence theorem inside a specified differential
field. The obstruction survives an infinitesimal coefficient:
$Dy=i\omega^{-1}y$ also has no nonzero solution. In contrast,
$Dy=i\omega^{-2}y$ has a solution given by a legitimate infinitesimal Hahn
exponential.

Theorem~\ref{thm:phase-criterion} explains all three examples at once.
Subsequent sections distinguish existence in the full class $\K$ from
existence in a fixed Hahn field, and then distinguish both from solving an
ordinary differential equation in an external analytic variable.

\subsection{What is imported and what is proved}

The deep input is the existence and surjectivity of the chosen surreal
derivation. We do not reproduce its construction through log-atomic numbers,
path derivatives and nested truncation ranks. We also use the established
normal-form description of $\No$, the real surreal exponential, and the
standard Hahn summability lemma. A later section records a restricted
composition theorem for omega-series as an additional, explicitly separate
input.

All phase criteria, exact-sequence assertions, workspace image calculations,
constant-coefficient classifications, and the displayed Picard--Vessiot
example are proved below from their stated inputs. They are presented as
expository deductions, without a novelty or priority claim. The accompanying
program verifies finite algebraic examples, not the imported structure theory
or the general theorems.

\section{Scalars, supports, and the chosen derivation}
\label{sec:setup}

\subsection{Set-sized supports in a class-sized field}

Write $\No$ for the surreal ordered field and $\K=\No[i]$ for its algebraic
closure. Every individual normal form has a set-sized support. It is
convenient to orient exponents by a formal infinitesimal symbol:
\begin{equation}\label{eq:normal}
 t^\gamma:=\omega^{-\gamma},\qquad
 x=\sum_{\gamma\in S}c_\gamma t^\gamma,
 \qquad c_\gamma\in\C,
\end{equation}
where $S\subseteq\No$ is a well-ordered set. For real $x$, the coefficients
are real. The least exponent of a nonzero normal form is its valuation
$v(x)$. Set $v(0)=+\infty$ as a notational convention.

In \eqref{eq:normal}, $\gamma\mapsto\omega^{-\gamma}$ is the Conway monomial
map. For arbitrary surreal exponents it must not be silently replaced by
$\exp(-\gamma\log\omega)$. The rational powers used in our explicit
workspace do agree with the usual field powers, which is all the derivative
calculation there requires.

We use definable classes, or a class-theoretic presentation with enough
replacement for the stated constructions. No support or family being summed
is a proper class. Statements about $\K$ mean elementwise class statements;
ordinary field constructions are applied to set-sized data. In particular,
our exact sequence in Section~\ref{sec:obstruction} is not a quotient set
whose members are proper classes.

\subsection{Finite, infinitesimal, and purely infinite parts}

Define
\[
 \OO=\{x\in\No:\abs{x}<n\text{ for some }n\in\N\},\qquad
 \mm=\{x\in\No:\abs{x}<1/n\text{ for every }n\geq1\}.
\]
The finite elements form a convex valuation ring and $\mm$ is its maximal
ideal. For $x\in\OO$, the standard part $\st(x)\in\R$ is the unique real
number such that $x-\st(x)\in\mm$. For complexified scalars put
$\OO_\K=\OO+i\OO$ and $\mm_\K=\mm+i\mm$.

Let $\PP$ be the additive class of purely infinite normal forms, including
zero: in the convention \eqref{eq:normal}, only negative exponents occur.
Every real surreal has the unique decomposition
\begin{equation}\label{eq:three-parts}
 x=\pipart(x)+\ct(x)+\epsilon,
 \qquad \pipart(x)\in\PP,\quad \ct(x)\in\R,\quad \epsilon\in\mm.
\end{equation}
Here $\ct(x)=[t^0]x$ is defined for \emph{every} surreal; it is standard part
only when $x$ is finite. Set $\fin(x)=\ct(x)+\epsilon$. The three projections
are additive and real-linear. The analogous coefficient projection on
$\K$ is complex-linear. They are not all field homomorphisms; for example,
$\ct(\omega)=\ct(t)=0$ but $\ct(\omega t)=1$.

\subsection{Hahn summability}

\begin{definition}
A set-indexed family $(x_j)_{j\in J}$ of Hahn normal forms is
\emph{summable} when the union of its supports is well ordered and every
exponent belongs to the support of only finitely many $x_j$. Its sum is
formed coefficientwise.
\end{definition}

We use the following support theorem in its additive-exponent form
\cite{Neumann}. If $S$ is a well-ordered set of strictly positive elements
of an ordered abelian group, its generated additive monoid $S^*$ is well
ordered, and each exponent has only finitely many presentations as a finite
sum of elements of $S$, up to ordering. This guarantees summability of power
series in an infinitesimal, and justifies the rearrangements of those series
used below. The statement is about supports, not numerical convergence.

In the full fine topology on $\No$ or $\K$, set-indexed convergence is highly
degenerate; the repository already explains this. None of our Hahn identities
is asserted to be a limit of partial sums in that topology. Conversely,
set-sized Hahn fields can have nontrivial valuation-topological convergence;
the full-class topological statements must not simply be transferred to
them \cite{RepoAnalysis,RepoFound}.

\subsection{The structural input}

\begin{theorem}[Berarducci--Mantova, imported]\label{thm:BM}
There is a specified surreal derivation $D:\No\to\No$ with
\[
 D(xy)=xDy+yDx,\qquad
 D\!\left(\sum_jx_j\right)=\sum_jDx_j,
 \qquad D(\exp x)=\exp(x)Dx.
\]
The middle assertion includes summability of the derivative family whenever
the original family is summable. Moreover,
\[
 \ker D=\R,\qquad D\omega=1,\qquad D(\No)=\No,
 \qquad x>\R\ \Longrightarrow\ Dx>0.
\]
\end{theorem}

This is the normalized ``simplest'' derivation of
\cite[Theorems 6.30 and 7.7]{BM}. The displayed properties are a usable
interface, not a new construction or a uniqueness characterization. There
are other surreal derivations. All intrinsic statements in this article
refer to this fixed $D$.

The real exponential is the usual Gonshor surreal exponential, extending
real exponentiation; its inverse on positive surreals is $\log$. Rational
powers and positive square roots are the ordinary field operations. These
background operations and normal forms are described in \cite{Gonshor}.

\section{Elementary intrinsic calculus}
\label{sec:real-calculus}

A field derivation is an additive map satisfying Leibniz' rule. Its input
is a scalar, not a function. Thus $D\omega=1$ is a statement about the chosen
structure on the field, whereas differentiation of the constant function
$z\mapsto\omega$ with respect to $z$ gives zero.

\subsection{Quotients, powers, and real exponentials}

Leibniz' rule gives $D1=0$, and hence, for $x\ne0$,
\begin{equation}\label{eq:quotient}
 D(x^{-1})=-x^{-2}Dx,
 \qquad
 D(x/y)=\frac{yDx-xDy}{y^2}.
\end{equation}
If $x>0$ and $q\in\Q$, differentiating a suitable equation for a rational
power gives
\begin{equation}\label{eq:rational-deriv}
 D(x^q)=q x^{q-1}Dx.
\end{equation}
Applying $D$ to $\exp(\log x)=x$ yields
\begin{equation}\label{eq:log-deriv}
 D\log x=\frac{Dx}{x}.
\end{equation}
In particular, with $t=\omega^{-1}$,
\[
 Dt=-t^2,\qquad D\log\omega=t,\qquad
 D(t^q)=-q t^{q+1}\quad(q\in\Q).
\]

Every real coefficient $a\in\No$ is a logarithmic derivative: choose $A$
with $DA=a$, then $D(e^A)/e^A=a$. Notice the two distinct uses of structure:
surjectivity gives $A$, and compatibility with the real exponential gives
$e^A$. Merely knowing $D$ is onto would not justify the second step.

\subsection{A useful consequence of the order axiom}

\begin{lemma}[Bounded derivatives are infinitesimal]\label{lem:bounded}
One has $D(\OO)\subseteq\mm$ and $D(\OO)=D(\mm)$.
\end{lemma}
\begin{proof}
Fix $f\in\OO$ and a positive integer $n$. Both $\omega+nf$ and
$\omega-nf$ are positive infinite, so the last axiom of
Theorem~\ref{thm:BM} gives
\[
 1+nDf>0,\qquad 1-nDf>0.
\]
Thus $\abs{Df}<1/n$ for every positive integer $n$. This says $Df\in\mm$.
Finally $f=\st(f)+\epsilon$ with $\epsilon\in\mm$, and the real standard
part is a constant of $D$. Therefore $Df=D\epsilon$.
\end{proof}

The conclusion is stronger than mere preservation of the finite valuation
ring. It will supply the elementary obstruction to intrinsic oscillation.
It is not the assertion $D(\mm)=\mm$; that equality is false. In fact
$t\notin D(\mm)$, as the phase criterion and $D\log\omega=t$ will show.

\subsection{Differentiation does not choose an independent variable}

One may informally think of $D$ as differentiation ``with respect to
$\omega$'' on simple expressions such as $\omega^3$ or $\log\omega$.
This is a useful mnemonic, but not a definition on all of $\No$. For a
general normal form, its exponents can themselves have nontrivial
intrinsic derivatives, and its monomials are not generic exponential powers.
An ordinary symbolic power rule applied indiscriminately to
$\omega^\gamma$ would not construct the Berarducci--Mantova derivation.

The distinction matters computationally. The formula
$D(t^q)=-q t^{q+1}$ is an exact rule for rational $q$, and extends strongly
to our rational Hahn workspace. It is not a general formula for surreal
$\gamma$ with $q$ replaced by $\gamma$.

\section{Complexification and infinitesimal phase}
\label{sec:phase}

\subsection{The unique extension to surcomplex scalars}

\begin{proposition}\label{prop:complexify}
The unique extension of $D$ to $\K=\No[i]$ is
\[
 D(a+ib)=Da+iDb.
\]
It is strongly additive, commutes with conjugation, has kernel $\C$, and
is surjective.
\end{proposition}
\begin{proof}
Any extension must satisfy $0=D(i^2+1)=2iDi$, so $Di=0$.
The displayed formula is then forced. Direct expansion proves Leibniz'
rule. The two real components of a summable complex family are summable;
strong additivity follows from the real statement. The kernel is
$\R+i\R=\C$. Given $a+ib$, choose real primitives $A,B$; then
$D(A+iB)=a+ib$. Conjugation commutes with the displayed formula.
\end{proof}

For $z\ne0$, set $r=\abs z=\sqrt{z\bar z}>0$ and $u=z/r$. Then
$u\bar u=1$. Differentiating $r^2=z\bar z$ proves
\begin{equation}\label{eq:modulus-deriv}
 \frac{Dr}{r}=\Ree\left(\frac{Dz}{z}\right).
\end{equation}
This is an algebraic identity over a real closed field with a derivation;
there is no contour, path, or topological limit involved.

\subsection{The local exponential and logarithm}

For $h\in\mm_\K$ define the Hahn sums
\begin{equation}\label{eq:local-functions}
 \EE(h)=\sum_{n=0}^{\infty}\frac{h^n}{n!},\qquad
 \LL(1+h)=\sum_{n=1}^{\infty}\frac{(-1)^{n+1}h^n}{n}.
\end{equation}
If $h=0$, these have their usual values. Otherwise the support of $h$ is
positive and well ordered, so the support theorem from
Section~\ref{sec:setup} applies. Every sum and rearrangement in the following
proof takes place over a set-sized positive generated monoid.

\begin{proposition}[Local exponential calculus]\label{prop:local-exp}
The maps
\[
 \EE:(\mm_\K,+)\longrightarrow(1+\mm_\K,\cdot),\qquad
 \LL:1+\mm_\K\longrightarrow\mm_\K
\]
are inverse group isomorphisms. They commute with conjugation, and
\begin{equation}\label{eq:Dlocal}
 D\EE(h)=\EE(h)Dh,\qquad
 D\LL(1+h)=\frac{Dh}{1+h}.
\end{equation}
\end{proposition}
\begin{proof}
The binomial formula and finite coefficient extraction give
$\EE(h+k)=\EE(h)\EE(k)$. The formal characteristic-zero identities
$\log(\exp X)=X$ and $\exp(\log(1+X))=1+X$ can be substituted because
every coefficient calculation after substitution has only finitely many
contributors at a fixed Hahn exponent. They yield mutual inverseness.
Conjugation commutes with each finite coefficient calculation.

Strong additivity permits termwise differentiation, including the assertion
that the differentiated terms form a summable family. Leibniz' rule gives
$D(h^n)=n h^{n-1}Dh$. Thus
\[
 D\EE(h)=Dh\sum_{n\geq1}\frac{h^{n-1}}{(n-1)!}=Dh\,\EE(h).
\]
Similarly,
\[
 D\LL(1+h)=Dh\sum_{n\geq1}(-1)^{n+1}h^{n-1}
           =\frac{Dh}{1+h}.
\]
Multiplication by the single scalar $Dh$ preserves the relevant summable
families. These are Hahn identities, not interchange-of-limit arguments.
\end{proof}

\subsection{Every unit direction has an infinitesimal phase defect}

Let $\Torus(\No)=\{u\in\K:u\bar u=1\}$. If $u=a+ib$ lies on this
circle, then $a^2+b^2=1$, so both coordinates are finite. Its standard part
$c=\st(u)$ belongs to the ordinary complex unit circle.

\begin{theorem}[Local phase decomposition]\label{thm:unit-phase}
Every $u\in\Torus(\No)$ has a unique expression
\begin{equation}\label{eq:unit-phase}
 u=c\,\EE(i\epsilon),\qquad
 c\in\C,\quad\abs c=1,\quad\epsilon\in\mm.
\end{equation}
Consequently
\begin{equation}\label{eq:unit-logder}
 \frac{Du}{u}=iD\epsilon.
\end{equation}
\end{theorem}
\begin{proof}
Take $c=\st(u)$ and $w=u/c\in1+\mm_\K$. The relation $w\bar w=1$
gives $\bar w=w^{-1}$. Proposition~\ref{prop:local-exp} implies
\[
 \overline{\LL(w)}=\LL(\bar w)=\LL(w^{-1})=-\LL(w).
\]
Hence $\LL(w)=i\epsilon$ for a real infinitesimal $\epsilon$, and
$w=\EE(i\epsilon)$. Conversely every expression in \eqref{eq:unit-phase}
has norm one. Its standard part recovers $c$, and the injective local
logarithm recovers $\epsilon$. Differentiating, and using $Dc=0$, proves
\eqref{eq:unit-logder}.
\end{proof}

This proof does not choose a global value of $\sin\omega$, a branch of a
global logarithm, or a global circular character. The ordinary phase $c$
is a differential constant. All variable intrinsic phase of a unit is
contained in its infinitesimal defect $\epsilon$.

\subsection{The simplest no-oscillation theorem}

\begin{proposition}\label{prop:no-oscillation}
The equations $Dy=iy$ and $Dy=-iy$ have only the zero solution in $\K$.
\end{proposition}
\begin{proof}
Suppose $Dy=iy$ and $y\ne0$. Conjugation gives $D\bar y=-i\bar y$, so
$D(y\bar y)=0$. Since $y\bar y>0$, it is a positive real constant.
Writing $y=a+ib$, both $a$ and $b$ are bounded. Their derivatives are
infinitesimal by Lemma~\ref{lem:bounded}. But the equation says
$Da=-b$ and $Db=a$, so both $a$ and $b$ are infinitesimal. Then
$a^2+b^2$ is infinitesimal, contradicting its being a positive real number.
Conjugation proves the assertion with $-i$.
\end{proof}

The argument makes the obstruction visible before any general criterion is
introduced: a rotating pair would have constant nonzero real norm, but
bounded surreal coordinates can only have infinitesimal derivatives.

\section{An exact criterion for first-order equations}
\label{sec:first-order}

\subsection{Finite accumulated phase}

For a real coefficient $b\in\No$, a \emph{phase primitive} means a scalar
$B\in\No$ satisfying $DB=b$. It is an intrinsic antiderivative, not an
integral along a time interval. Surjectivity guarantees its existence, and
any two choices differ by a real constant. Thus whether $B$ is finite is
independent of the choice.

\begin{theorem}[First-order phase criterion]\label{thm:phase-criterion}
Let $q=a+ib\in\K$, with $a,b\in\No$. Choose $A,B\in\No$ such that
$DA=a$ and $DB=b$. The following conditions are equivalent:
\begin{enumerate}[label=(\roman*)]
\item $Dy=qy$ has a nonzero solution in $\K$;
\item $b\in D(\mm)$;
\item $B\in\OO$;
\item $\pipart(B)=0$.
\end{enumerate}
When these conditions hold, put $\epsilon=B-\st(B)$. All solutions are
\begin{equation}\label{eq:all-first-order}
 y=c\,e^A\EE(i\epsilon),\qquad c\in\C.
\end{equation}
\end{theorem}
\begin{proof}
Suppose $y\ne0$ solves the equation. Write $y=ru$, where $r=\abs y>0$
and $u\bar u=1$. Equation~\eqref{eq:modulus-deriv} gives $Dr/r=a$.
Theorem~\ref{thm:unit-phase} writes $u=c_0\EE(i\epsilon)$ with
$\epsilon\in\mm$. Taking the imaginary part of
\[
 \frac{Dy}{y}=\frac{Dr}{r}+\frac{Du}{u}
\]
gives $b=D\epsilon$. This proves (i)$\Rightarrow$(ii).

If $b=D\epsilon$ with $\epsilon\in\mm$, then
$D(B-\epsilon)=0$. Hence $B-\epsilon\in\R$ and $B$ is finite.
Conversely, a finite $B$ differs from an infinitesimal by its real standard
part, so $b=D(B-\st(B))$. This proves (ii)$\Leftrightarrow$(iii), and
\eqref{eq:three-parts} gives (iii)$\Leftrightarrow$(iv).

If (iii) holds, the nonzero scalar $y_0=e^A\EE(i\epsilon)$ satisfies
$Dy_0/y_0=a+iD\epsilon=q$. This proves existence. For any other solution,
\[
 D(y/y_0)=\frac{y_0Dy-yDy_0}{y_0^2}=0,
\]
so $y/y_0\in\C$. This includes the zero solution and proves
\eqref{eq:all-first-order}.
\end{proof}

The criterion involves a primitive of the imaginary coefficient, not just
its size. In particular,
\begin{equation}\label{eq:strict-small}
 D(\mm)\subsetneq\mm.
\end{equation}
Indeed $D\log\omega=t$, while $\log\omega$ is infinite, so
$t\notin D(\mm)$. A coefficient can be infinitesimal and still accumulate
an infinite intrinsic phase.

\subsection{Examples across two scales}

All exponentials $\EE$ in the following table have infinitesimal arguments
and therefore denote the support-certified sums of
\eqref{eq:local-functions}. Set $L=\log\omega$.

\begin{center}
\small
\begin{tabular}{@{}p{0.22\textwidth}p{0.25\textwidth}p{0.42\textwidth}@{}}
\toprule
Imaginary coefficient $b$ & A primitive $B$ & Nonzero solution of $Dy=iby$?\\
\midrule
$1$ & $\omega$ & No: the phase primitive is infinite.\\[4pt]
$t$ & $L$ & No, although $b$ is infinitesimal.\\[4pt]
$t^2$ & $-t$ & Yes: $\EE(-it)$.\\[4pt]
$t^{3/2}$ & $-2t^{1/2}$ & Yes: $\EE(-2it^{1/2})$.\\[4pt]
$1/(\omega L)$ & $\log L$ & No: a slower divergence still obstructs.\\[4pt]
$1/(\omega L^2)$ & $-1/L$ & Yes: $\EE(-i/L)$.\\
\bottomrule
\end{tabular}
\end{center}

Each derivative is checked directly from \eqref{eq:quotient} and
\eqref{eq:log-deriv}. The familiar growth facts
$1\ll\log\omega\ll\omega^\delta$ for rational $\delta>0$, and the
corresponding fact that $\log\log\omega$ is infinite, are properties of the
surreal exponential and logarithm \cite{Gonshor}. The table does not use a
notion of integration to a numerical endpoint.

Adding any real coefficient $a$ introduces no further obstruction in the
full class: multiply the displayed solution by $e^A$, where $DA=a$.
That assertion can fail in a restricted workspace because $A$ or $e^A$
may leave the workspace. Section~\ref{sec:hahn} computes that failure
exactly in a basic example.

\subsection{What the result does not say}

Theorem~\ref{thm:phase-criterion} is a complete classification for scalar
homogeneous equations of order one over the chosen differential class field.
It is not a claim that all nonlinear differential equations have solutions,
that the field is differentially closed, or that a primitive can be computed
from every finite input description. In fact Proposition~\ref{prop:no-oscillation}
already disproves differential closedness: the compatible conditions
$DY=iY$ and $Y\ne0$ acquire a solution in a differential extension, as
Section~\ref{sec:PV} constructs explicitly, but not in $\K$.

\section{The obstruction map and the maximal exponential domain}
\label{sec:obstruction}

\subsection{A concrete replacement for a class quotient}

Define the logarithmic derivative of a nonzero scalar by
\[
 \Dlog:\K^\times\longrightarrow(\K,+),\qquad
 \Dlog(y)=\frac{Dy}{y}.
\]
It is a group homomorphism because $\Dlog(yz)=\Dlog(y)+\Dlog(z)$.
Its kernel is $\C^\times$.

For $q=a+ib$, choose a real phase primitive $B$ and define
\begin{equation}\label{eq:obstruction}
 \Omega(q)=\pipart(B)\in\PP.
\end{equation}
No global selection of unrelated primitives is necessary for this
definition: any two primitives differ by a real constant, which has zero
purely infinite part.

\begin{theorem}[Logarithmic-derivative obstruction]\label{thm:exact}
The map $\Omega$ is a well-defined surjective real-linear map, and
\begin{equation}\label{eq:ld-image}
 \im\Dlog=\No+iD(\mm)=\ker\Omega.
\end{equation}
There is an elementwise exact sequence of class groups
\begin{equation}\label{eq:exact}
 1\longrightarrow\C^\times\longrightarrow\K^\times
 \xrightarrow{\ \Dlog\ }(\K,+)
 \xrightarrow{\ \Omega\ }(\PP,+)\longrightarrow0.
\end{equation}
\end{theorem}
\begin{proof}
If $B_1$ and $B_2$ are primitives of $b_1$ and $b_2$, then
$B_1+B_2$ is a primitive of $b_1+b_2$. Purely infinite projection is
additive; multiplication by a real constant works in the same way. Thus
$\Omega$ is real-linear. For any $P\in\PP$,
\[
 \Omega(iDP)=P,
\]
so it is onto. Theorem~\ref{thm:phase-criterion} says precisely that
$q$ lies in the logarithmic-derivative image if and only if
$\Omega(q)=0$, and identifies that image with
$\No+iD(\mm)$. The remaining exactness assertions follow from the kernel
calculation and the inclusion of ordinary nonzero complex constants.
\end{proof}

Writing the target as $\PP$ avoids treating a coset of $\OO$ in $\No$ as
an element of a set: such a coset is generally a proper class. Informally,
$\PP$ realizes the additive quotient ``$\No/\OO$'' through canonical
normal-form representatives. Equation~\eqref{eq:exact} means the stated
kernel and image identities, not the formation of a set of all such cosets.

The image of $\Dlog$ is real-linear but is not complex-linear. For example,
$1=\Dlog(e^\omega)$ lies in the image, while $i$ does not. An implementation
must not model this image as a complex vector subspace.

\subsection{An intrinsic exponential has a sharp domain restriction}

An exponential-compatible value at an argument $z$ would be a nonzero scalar
$E(z)$ with
\begin{equation}\label{eq:intrinsic-exp-law}
 D(E(z))=E(z)Dz.
\end{equation}
This is not external differentiation of $E$ as a function. It is a condition
on the result of applying $D$ to the scalar $E(z)$.

\begin{theorem}[Maximal pointwise exponential domain]\label{thm:strip}
For $z=a+ib\in\K$, a nonzero value satisfying
\eqref{eq:intrinsic-exp-law} exists if and only if $b$ is finite.
On the additive subgroup
\[
 \PS=\No+i\OO
\]
there is a group homomorphism
\begin{equation}\label{eq:strip-exp}
 E_{\PS}(a+ib)=e^a e^{i\st(b)}
              \EE\bigl(i(b-\st(b))\bigr)
\end{equation}
that satisfies \eqref{eq:intrinsic-exp-law} and extends the ordinary
complex exponential. No extension to all of $\K$ can satisfy that law.
\end{theorem}
\begin{proof}
Apply Theorem~\ref{thm:phase-criterion} to the coefficient
$Dz=Da+iDb$. A primitive of its imaginary part is $b$ itself, so existence
is equivalent to $b\in\OO$.

For such $b$, all factors in \eqref{eq:strip-exp} are defined. Additivity of
standard part on finite scalars, the real exponential law, and the local
exponential group law prove multiplicativity with respect to addition of
arguments. The ordinary factor $e^{i\st(b)}$ is a differential constant.
Thus its intrinsic logarithmic derivative is
\[
 Da+iD(b-\st(b))=Da+iDb=Dz.
\]
At $z=i\omega$, the required equation would be $DY=iY$ with $Y\ne0$,
which is impossible.
\end{proof}

The maximality here is stronger than a failure of one candidate formula:
there is no admissible nonzero value at an excluded point. On the other
hand, we do not claim that the differential law by itself uniquely specifies
every normalization on $\PS$; ratios of admissible values are differential
constants. Formula~\eqref{eq:strip-exp} is the specified canonical local
construction with the ordinary complex normalization.

The domain $\PS$ is $D$-stable because $D\OO\subseteq\mm$, but it is not
a subring. For example, $i\in\PS$ and $\omega\in\PS$, while
$i\omega\notin\PS$. A symbolic system should therefore attach a domain
predicate to $E_{\PS}$ rather than advertise it as an exponential operation
on a differential subfield.

\subsection{Why this does not contradict global phase conventions}

The repository's trigonometric report studies finite phases and possible
global extensions of circular group laws, and explicitly does not select
an intrinsic derivation \cite{RepoTrig}. Theorem~\ref{thm:strip} does not
invalidate those algebraic constructions. It says that any assignment of a
global phase must fail the additional law $D(e^{ib})=ib'e^{ib}$ at some
infinite arguments when $D$ is the derivation fixed here. At $b=\omega$,
that incompatibility is unavoidable and independent of which phase was
chosen.

Likewise the ordinary function $z\mapsto e^{iz}$ satisfies an external
analytic differential equation. Its ordinary-variable derivative cannot
be substituted for $D$ acting on a supposed value at $\omega$. The bridge
between these operations requires extra structure; it is not supplied by
using the same symbol for both.

\section{Primitives, normalization, and inhomogeneous equations}
\label{sec:primitives}

\subsection{A canonical linear primitive operator}

Surjectivity and the known constant field give a convenient exact
normalization.

\begin{proposition}\label{prop:J}
There is a unique complex-linear map $J:\K\to\K$ such that
\begin{equation}\label{eq:J}
 DJf=f,\qquad \ct(Jf)=0.
\end{equation}
It also satisfies
\begin{equation}\label{eq:JD}
 JDf=f-\ct(f).
\end{equation}
\end{proposition}
\begin{proof}
Take any primitive $F$ of $f$ and set $Jf=F-\ct(F)$. Any other primitive
is $F+c$ with $c\in\C$, so the result is independent of that choice.
The normalized primitive is unique because its difference from any other
normalized primitive is a complex constant with zero constant coefficient.
Additivity and complex-linearity follow by uniqueness. The scalar
$f-\ct(f)$ is the normalized primitive of $Df$, proving \eqref{eq:JD}.
\end{proof}

This is a semantic construction by a unique characterization, not a
terminating algorithm for arbitrary inputs. It also does not require a
separate axiom of global choice: a unique primitive exists after the
normalization has been specified.

\subsection{The integration constant is not evaluation at an endpoint}

Integration by parts has the exact algebraic form
\begin{equation}\label{eq:parts}
 J(fDg+gDf)=fg-\ct(fg).
\end{equation}
The final term is not generally $\ct(f)\ct(g)$. This changes a familiar
algebraic identity for a normalized integration operator.

\begin{proposition}\label{prop:RB}
For $f,g\in\K$,
\begin{equation}\label{eq:RB-correction}
 Jf\,Jg
 =J(fJg+Jf\,g)+\ct(Jf\,Jg).
\end{equation}
The correction term cannot in general be omitted.
\end{proposition}
\begin{proof}
Differentiate $Jf\,Jg$ and apply \eqref{eq:JD}. This gives the identity.
For the final assertion take $f=1$ and $g=t^2$. Then
\[
 Jf=\omega,\qquad Jg=-t,\qquad Jf\,Jg=-1,
 \qquad fJg+Jf\,g=-t+\omega t^2=0.
\]
The version without the correction would assert $-1=0$.
\end{proof}

In terminology sometimes used for integral operators, $J$ is not a
weight-zero Rota--Baxter operator on all of $\K$. Equation~\eqref{eq:RB-correction}
is the precise replacement. No such terminology is needed to use the
formula: the important point is that the selected constant coefficient is
linear but not multiplicative. A ``zero constant of integration'' convention
is therefore not automatically an integral based at a common evaluation
point.

\subsection{Variation of constants with its necessary hypothesis}

\begin{theorem}\label{thm:inhom}
Suppose $y_0\ne0$ satisfies $Dy_0=qy_0$ in $\K$. For every $r\in\K$, all
solutions of $Dy=qy+r$ are
\begin{equation}\label{eq:variation}
 y=y_0\bigl(J(r/y_0)+c\bigr),\qquad c\in\C.
\end{equation}
If the corresponding homogeneous equation has no nonzero solution, an
inhomogeneous equation has at most one solution in $\K$.
\end{theorem}
\begin{proof}
The quotient rule gives
$D(y/y_0)=(Dy-qy)/y_0$. Thus the equation is equivalent to
$D(y/y_0)=r/y_0$, and Proposition~\ref{prop:J} gives
\eqref{eq:variation}. In the second case, the difference of two solutions
would be a nonzero homogeneous solution unless they were equal.
\end{proof}

No conclusion about existence in the second case follows from variation of
constants, because its integrating factor does not exist in $\K$. This
article does not invoke the additional differential-field structure that
would be needed for a general inhomogeneous existence theorem. The absence
of a homogeneous solution also does not itself imply inhomogeneous
nonexistence. For example,
\[
 (D-i)i=1,\qquad (D-i)(i\omega+1)=\omega.
\]
These are unique solutions of their respective equations in $\K$.

For $P\in\C[\omega]$ of degree at most $N$, a finite inverse formula is
\begin{equation}\label{eq:poly-inverse}
 (D-i)^{-1}P
 =i\sum_{k=0}^{N}(-i)^kD^kP.
\end{equation}
Indeed $D-i=-i(1+iD)$ and $D^{N+1}P=0$, so the finite geometric identity
applies. It is a true terminating formula, with no infinite summation or
asymptotic truncation.

\section{Constant-coefficient differential equations}
\label{sec:constant-ODE}

The phase obstruction controls more than a single first-order example.
It completely determines the solution space of a scalar homogeneous
constant-coefficient equation.

\subsection{Repeated zero eigenvalues}

\begin{lemma}\label{lem:Dpowers}
For every positive integer $m$,
\[
 \ker D^m=\C[\omega]_{<m},
\]
where the right-hand side consists of polynomials in $\omega$ of degree
strictly less than $m$ with ordinary complex coefficients.
\end{lemma}
\begin{proof}
The inclusion from right to left follows from $D\omega=1$.
For the reverse inclusion induct on $m$. For $m=1$, this is the constant
field statement. If $D^mf=0$, then $D^{m-1}f=c\in\C$. Subtract
$c\omega^{m-1}/(m-1)!$ from $f$; the result lies in $\ker D^{m-1}$ and
hence is a polynomial of degree at most $m-2$.
\end{proof}

For a real constant $\lambda$, multiplication by $e^{\lambda\omega}$
conjugates $D$ to $D-\lambda$:
\begin{equation}\label{eq:conjugate}
 (D-\lambda)(e^{\lambda\omega}f)=e^{\lambda\omega}Df.
\end{equation}
For a nonreal ordinary constant $\lambda=a+ib$, with $b\ne0$, the phase
primitive $b\omega$ is infinite, so $D-\lambda$ is injective. Consequently
all its positive powers are injective as well.

\subsection{The classification theorem}

\begin{theorem}\label{thm:constant-ODE}
Let $P(X)\in\C[X]$ be nonzero. For each real root $\lambda$ of $P$, let
$m_\lambda$ be its multiplicity. Then
\begin{equation}\label{eq:constant-solutions}
 \ker P(D)=
 \bigoplus_{\substack{\lambda\in\R\\P(\lambda)=0}}
 e^{\lambda\omega}\C[\omega]_{<m_\lambda}.
\end{equation}
In particular, its dimension over $\C$ is the total multiplicity of the
\emph{real} roots of $P$, not necessarily the degree of $P$.
\end{theorem}
\begin{proof}
If $P$ is a nonzero constant, both sides of the claimed identity are zero.
Otherwise factor $P$ into its pairwise coprime primary factors
$p_j(X)=(X-\lambda_j)^{m_j}$. Multiplication by a nonzero scalar factor
in $P$ does not change the kernel. For pairwise coprime factors there are
polynomials $e_j$ satisfying
\[
 \sum_j e_j(X)=1,\qquad
 e_j\equiv1\pmod{p_j},\qquad
 e_j\equiv0\pmod{p_k}\quad(k\ne j).
\]
These are the usual Bezout projectors modulo $P$; their sum may first be
$1$ modulo $P$, which is all that is needed after application to $\ker P(D)$.
If $P(D)y=0$, put $y_j=e_j(D)y$. The product $p_j e_j$ is divisible by
$P$, so $p_j(D)y_j=0$, and $y=\sum_jy_j$. On each primary kernel the
corresponding projector is identity and the others are zero, so the sum is
direct.

For real $\lambda_j$, repeated use of \eqref{eq:conjugate} and
Lemma~\ref{lem:Dpowers} gives
\[
 \ker(D-\lambda_j)^{m_j}
 =e^{\lambda_j\omega}\C[\omega]_{<m_j}.
\]
For nonreal $\lambda_j$, that primary kernel is zero by injectivity.
Combining the primary kernels proves \eqref{eq:constant-solutions}.
The indicated polynomial spaces have the asserted finite dimensions.
\end{proof}

\begin{example}
The equation $D^2y+y=0$ has only the zero solution, since the roots of
$X^2+1$ are nonreal. The equation
\[
 (D-2)^3(D+1)(D^2+1)y=0
\]
has all its solutions in the four-dimensional space
\[
 y=e^{2\omega}(c_0+c_1\omega+c_2\omega^2)+c_3e^{-\omega},
 \qquad c_j\in\C.
\]
The two formal oscillatory modes belonging to $\pm i$ are absent.
\end{example}

\subsection{Constant matrices}

Let $A\in M_n(\C)$. An ordinary constant change of basis puts $A$ in
Jordan form without changing how $D$ acts on vectors. In a block
$\lambda I+N$ of size $m$, the equations read, with the superdiagonal
convention,
\[
 (D-\lambda)y_j=y_{j+1}\ (j<m),\qquad
 (D-\lambda)y_m=0.
\]
If $\lambda$ is nonreal, injectivity forces $y_m=0$, then $y_{m-1}=0$,
and so on. If $\lambda$ is real, a fundamental block is
\[
 e^{\lambda\omega}\exp(\omega N)
 =e^{\lambda\omega}\sum_{k=0}^{m-1}\frac{\omega^kN^k}{k!}.
\]
The sum is finite because $N^m=0$.

\begin{corollary}\label{cor:matrix}
The system $DY=AY$ with $A\in M_n(\C)$ has a fundamental matrix in
$\operatorname{GL}_n(\K)$ if and only if every eigenvalue of $A$ is real.
Its vector-solution space always has dimension equal to the total algebraic
multiplicity of the real eigenvalues.
\end{corollary}

This theorem concerns ordinary constant matrices, not arbitrary
surcomplex-coefficient systems. It is already sufficient to rule out a
universal matrix-exponential constructor in the intrinsic field: the
rotation matrix
$\left(\begin{smallmatrix}0&-1\\1&0\end{smallmatrix}\right)$
has no nonzero solution vector, while a real diagonal matrix does have the
expected exponential modes.

\section{A completely explicit Hahn workspace}
\label{sec:hahn}

The full-class results above do not automatically describe a symbolic
backend or a set-sized coefficient domain. A particularly informative
workspace is
\begin{equation}\label{eq:HahnQ}
 \HH=\C((t^\Q)),\qquad t=\omega^{-1},
\end{equation}
embedded in $\K$ through its normal forms. Let $\HH_\R=\R((t^\Q))$.
Every element of $\HH$ has a well-ordered rational support; this is a full
Hahn field, not just the union of finite-ramification Laurent-series fields.
For example, well-ordered supports need not have a common denominator.

\subsection{The exact derivative formula}

\begin{proposition}\label{prop:Hahn-deriv}
The field $\HH$ is $D$-stable, and
\begin{equation}\label{eq:Hahn-D}
 D\left(\sum_{q\in S}c_qt^q\right)
 =\sum_{q\in S}(-q)c_qt^{q+1}.
\end{equation}
\end{proposition}
\begin{proof}
Use \eqref{eq:rational-deriv} on each monomial and strong additivity.
Translation by $1$ preserves well-ordering, and distinct exponents remain
distinct. The term with $q=0$ vanishes. Hence the right side is a valid
Hahn series with rational support.
\end{proof}

Thus this restriction of the abstract derivation is completely explicit.
This does not mean the full Berarducci--Mantova construction has been
implemented: the special rational monomial subgroup is doing essential
work.

\subsection{Additive primitives: one resonant coefficient}

\begin{theorem}[Derivative image in the rational Hahn field]
\label{thm:Hahn-integral}
For $f=\sum_q c_qt^q\in\HH$,
\begin{equation}\label{eq:Hahn-integrability}
 f\in D(\HH)\quad\Longleftrightarrow\quad [t^1]f=0.
\end{equation}
If $c_1=0$, its normalized primitive in $\HH$ is
\begin{equation}\label{eq:Hahn-integral}
 Jf=-\sum_{q\ne1}\frac{c_q}{q-1}t^{q-1}.
\end{equation}
For arbitrary $f\in\HH$, the normalized primitive in the ambient $\K$ is
\begin{equation}\label{eq:Hahn-log-integral}
 Jf=c_1\log\omega-
       \sum_{q\ne1}\frac{c_q}{q-1}t^{q-1}.
\end{equation}
\end{theorem}
\begin{proof}
In \eqref{eq:Hahn-D}, a contribution at exponent $1$ could only come from
$q=0$, whose coefficient is multiplied by zero. This proves necessity in
\eqref{eq:Hahn-integrability}. If $c_1=0$, translate the remaining support
by $-1$ and divide each coefficient by the nonzero rational number $q-1$.
The support stays well ordered, so \eqref{eq:Hahn-integral} defines an
actual element of $\HH$. Differentiation by \eqref{eq:Hahn-D} returns $f$.
There is no exponent-zero term because $q=1$ was excluded.

For general $f$, use $D\log\omega=t$. The surreal $\log\omega$ is purely
infinite, and hence has zero ordinary constant coefficient; therefore
\eqref{eq:Hahn-log-integral} is normalized as well.
\end{proof}

In particular, $\HH$ is closed under differentiation but not under
antidifferentiation. Since $t$ has no primitive in $\HH$, the known ambient
primitive $\log\omega$ does not belong to $\HH$. No appeal to topology,
cardinality, or computational limitations is needed for this conclusion.
The obstruction is the single resonant coefficient $[t]f$.

The formula remains valid for infinite rational supports. There is no need
to assume a denominator bound or a discrete sequence of increasing
exponents. Translation of an already well-ordered support is the precise
support certificate.

\subsection{Logarithmic derivatives inside the workspace}

A nonzero element $y\in\HH$ has a unique leading-term factorization
\[
 y=c\,t^r(1+h),\qquad c\in\C^\times,\quad r\in\Q,\quad v(h)>0,
\]
with $h=0$ permitted. The factor $1+h$ admits the local logarithm from
Proposition~\ref{prop:local-exp}.

\begin{theorem}[Logarithmic-derivative image in $\HH$]
\label{thm:Hahn-logder}
One has
\begin{equation}\label{eq:Hahn-logder-image}
 \Dlog(\HH^\times)
 =\{-r t+g:\ r\in\Q,\quad g=0\text{ or }v(g)>1\}.
\end{equation}
For $q=-r t+g$ in this image, put
\begin{equation}\label{eq:workspace-factor}
 k=-\sum_{\alpha>1}\frac{[t^\alpha]g}{\alpha-1}t^{\alpha-1}.
\end{equation}
Then all solutions in $\HH$ of $Dy=qy$ are
\begin{equation}\label{eq:Hahn-all}
 y=c\,t^r\EE(k),\qquad c\in\C.
\end{equation}
\end{theorem}
\begin{proof}
For the leading-term factorization above,
\[
 \frac{Dy}{y}=-r t+\frac{Dh}{1+h}
             =-r t+D\LL(1+h).
\]
The series $\LL(1+h)$ has positive valuation or is zero. Every nonzero
term in its derivative has exponent strictly greater than $1$, by
\eqref{eq:Hahn-D}. This proves containment in the right side of
\eqref{eq:Hahn-logder-image}.

Conversely, if $v(g)>1$, the series $k$ in
\eqref{eq:workspace-factor} has positive valuation and satisfies $Dk=g$.
The same conclusions hold for $g=k=0$. Thus
$\EE(k)\in\HH$ is defined by a positive-support Hahn sum, and
$\Dlog(t^r\EE(k))=-rt+g$. Finally, the ratio of any two nonzero
solutions is in the constant field $\C$.
\end{proof}

In this local construction, $\EE(k)-1$ has positive valuation or is zero;
there is no restriction to finite Hahn polynomials. The summability certificate comes from the positive
support of $k$.

The condition at exponent $1$ is much more restrictive for logarithmic
integration than for additive integration. A logarithmic derivative in
$\HH$ has no exponents below $1$, and its exponent-$1$ coefficient is a
real rational number. Its higher exponents may have arbitrary ordinary
complex coefficients.

\subsection{Ambient solutions versus workspace solutions}

\begin{proposition}\label{prop:Hahn-ambient}
Let $q=a+ib\in\HH$ with $a,b\in\HH_\R$. A nonzero solution of $Dy=qy$
exists in the ambient $\K$ if and only if
\begin{equation}\label{eq:Hahn-ambient}
 b=0\quad\text{or}\quad v(b)>1.
\end{equation}
\end{proposition}
\begin{proof}
Apply \eqref{eq:Hahn-log-integral} to $b$. If its support lies strictly
above $1$, the displayed primitive is infinitesimal.
If the least support exponent $\alpha$ is less than $1$, its primitive
has a nonzero leading term proportional to $t^{\alpha-1}$, which is
infinite. This term dominates all later rational-power terms and any
$\log\omega$ term. The growth comparison
$\log\omega\ll\omega^\delta$ for positive rational $\delta$ ensures the
last assertion. If the least exponent is $1$, the primitive is a nonzero
real multiple of $\log\omega$ plus an infinitesimal, hence is infinite.
Theorem~\ref{thm:phase-criterion} now gives exactly
\eqref{eq:Hahn-ambient}.
\end{proof}

Thus an ambient solution may exist but require an exponential or monomial
outside the rational Hahn workspace. These are three distinct outcomes,
not a binary ``solvable/unsolvable'' flag.

\begin{center}
\small
\begin{tabular}{@{}p{0.26\textwidth}p{0.16\textwidth}p{0.45\textwidth}@{}}
\toprule
Coefficient $q$ in $Dy=qy$ & Location & Explanation or one solution\\
\midrule
$1$ & $\K$, not $\HH$ & $e^\omega$; the coefficient at exponent $0$ violates \eqref{eq:Hahn-logder-image}.\\[4pt]
$\sqrt2\,t$ & $\K$, not $\HH$ & $\exp(\sqrt2\log\omega)$; the exponent-$1$ coefficient is irrational.\\[4pt]
$it$ & Neither & Its phase primitive is $\log\omega$.\\[4pt]
$it^2$ & $\HH$ & $\EE(-it)$.\\[4pt]
$-\tfrac32t+(1+2i)t^2$ & $\HH$ & $t^{3/2}\EE(-(1+2i)t)$.\\
\bottomrule
\end{tabular}
\end{center}

The first two rows mean no \emph{nonzero} solution belongs to $\HH$.
Multiplication by an ordinary complex constant cannot remove either
obstruction, because it does not change a logarithmic derivative.

\subsection{A resonant normal form}

There is a useful interpretation of Theorem~\ref{thm:Hahn-logder}.
A scalar differential equation over $\HH$ can have a fundamental solution
in the same field only when its logarithmic coefficient splits into a
rational monomial contribution $-rt$ and a higher-valuation perturbation.
The monomial supplies the exponent-$1$ residue, while the infinitesimal
exponential absorbs the perturbation. Lower exponents require larger-scale
real exponentials, and imaginary lower exponents face the stronger ambient
phase obstruction.

This resembles a one-dimensional normal-form calculation, but every
assertion here is algebraic and exact. There is no asymptotic error term in
\eqref{eq:Hahn-all}, and no claim that rational exponent $r$ is a numerical
approximation to a more general exponent.

\section{Differential localization with support control}
\label{sec:localization}

The repository's algebraic workspace principle lets a set of surcomplex
scalars be placed inside a set-sized Hahn field. A differential calculation
needs more: the workspace should contain the derivatives of its elements.
Closing only under field operations does not provide that condition.

\subsection{A derivative-stable Hahn workspace for every set of data}

For a set-sized divisible additive subgroup $\Gamma\subseteq\No$, write
\[
 K_\Gamma=\C((t^\Gamma))\subseteq\K.
\]
The coefficients are ordinary complex constants. This restriction is
important in the following proof.

\begin{theorem}[Differential Hahn localization]\label{thm:localization}
Every set $A\subseteq\K$ is contained in some set-sized full Hahn field
$K_\Gamma$ satisfying $D(K_\Gamma)\subseteq K_\Gamma$.
\end{theorem}
\begin{proof}
Let $\Gamma_0$ be the rational linear span of the union of the normal-form
supports of all members of $A$. This is a set. Recursively define
\begin{equation}\label{eq:Gamma-closure}
 \Gamma_{n+1}=\Span_\Q\left(
 \Gamma_n\cup\bigcup_{\gamma\in\Gamma_n}\supp D(t^\gamma)
 \right),\qquad
 \Gamma=\bigcup_{n\in\N}\Gamma_n.
\end{equation}
Replacement and union show that each stage and their union are sets.
Every stage is a rational vector space, so $\Gamma$ is divisible.

Fix $\gamma\in\Gamma$. It lies in some $\Gamma_n$, and therefore
$\supp D(t^\gamma)\subseteq\Gamma_{n+1}\subseteq\Gamma$.
Now take any
$f=\sum_{\gamma\in S}c_\gamma t^\gamma\in K_\Gamma$.
Strong additivity gives the summable family
$(c_\gamma D(t^\gamma))_{\gamma\in S}$. Each member has support in
$\Gamma$, so their sum also has support in $\Gamma$. Its coefficients
are in $\C$, and it equals $Df$. Thus $Df\in K_\Gamma$.
The starting choice of $\Gamma_0$ ensures $A\subseteq K_\Gamma$.
\end{proof}

An essential point is the final appeal to strong additivity. We did not
prove or use
\[
 K_{\bigcup_n\Gamma_n}=\bigcup_nK_{\Gamma_n},
\]
which is generally false: one Hahn series can use exponents appearing at
arbitrarily late stages. The correct proof handles every series over the
union simultaneously by its monomial supports.

This is a semantic existence construction, not an effective algorithm for
computing $D(t^\gamma)$ at arbitrary surreal exponents. It also supplies
no optimal birthday or cardinal bound. For more specialized full surreal
subfields stable under several transcendental and differential operations,
Bournez--Guilmant give a separate structural treatment \cite{BG}.

\subsection{Different closure promises are genuinely different}

The field $\HH$ in Section~\ref{sec:hahn} is already $D$-stable but does
not contain a primitive of $t$. Thus Theorem~\ref{thm:localization} cannot
be read as primitive closure. It does not assert closure under real
exponentiation, arbitrary logarithms, or solutions of all first-order
linear equations either.

For a set-sized \emph{subfield}, rather than a full Hahn field, one can
obtain a different useful result. Start with a field generated by the
specified data and the ordinary constants; at each finite stage adjoin
values of a specified set of finitary operations, such as $D$, $J$,
conjugation, and real exponential or logarithm where defined, then take the
field generated by those values. The union over the natural numbers is a
set-sized field stable under each listed operation. Any finite tuple from
the union occurs at a common stage, so its operation value occurs at the
next one.

This finitary closure proof is not a proof that the resulting field is a
full Hahn field, nor that it contains every summable family of its own
members. The distinction is useful for formalization: a small operational
workspace can support a particular proof even when it is not closed under
all possible Hahn sums.

\subsection{Changing the independent scale}

If $s\in\No$ has $Ds\ne0$, the rescaled operator
\[
 D_s f=\frac{Df}{Ds}
\]
is another derivation and satisfies $D_s s=1$. It need not retain the
chosen normalization at $\omega$ or all the same order conventions.
An equation $D_s y=qy$ is equivalent to $Dy=(Ds)q\,y$.
Consequently its phase test must be applied to the imaginary coefficient of
$(Ds)q$, not to that of $q$ alone. A backend cannot change the independent
scale while leaving differential solvability metadata untouched.

For example, take $s=\log\omega$, so $Ds=t$. Then $D_s y=iy$ is the
forbidden equation $Dy=ity$, whereas $D_sy=it y$ becomes
$Dy=it^2y$ and has a nonzero solution. The change of scale preserves the
underlying field, but it changes the coefficient that controls accumulated
phase.

\section{Four derivatives and an evaluation rule}
\label{sec:derivatives}

The safest way to connect this article to the existing analytic documents
is to give different operators different types. A formula can be correct
for one of them and false for another without any contradiction.

\subsection{Intrinsic, formal, coefficientwise, and fine derivatives}

The intrinsic $D$ acts on scalar values in $\K$ and has constants $\C$.
The formal derivative $\partial_Z$ on $\K[[Z]]$ differentiates the
indeterminate and treats every coefficient in $\K$ as a constant.
On a coherent Hahn family
\[
 F(z)=\sum_{\gamma\in S}f_\gamma(z)t^\gamma
\]
with each $f_\gamma$ holomorphic on one fixed ordinary domain $U\subseteq\C$,
the coefficientwise analytic derivative is
\[
 \partial_zF=\sum_{\gamma\in S}f'_\gamma(z)t^\gamma.
\]
Its support is a subset of the original support. Finally, a fine derivative
of a function $F$ at a scalar $x$ is an epsilon--delta assertion about
$[F(x+h)-F(x)]/h$ as the actual field variable $h$ becomes small. It requires
a function and a domain, not just a scalar field structure.

For example, in the coefficientwise family $F(z)=tz$ with ordinary $z$,
\begin{equation}\label{eq:deriv-separator}
 \partial_zF(z)=t,\qquad D(F(z))=-t^2z.
\end{equation}
The second equality holds because that fixed ordinary complex $z$ is a
constant of $D$. It is not differentiating the symbol $z$ as an independent
variable.

\subsection{Commuting derivations on formal series}

Define coefficient differentiation on $\K[[Z]]$ by
\[
 \DD\left(\sum_{n\geq0}a_nZ^n\right)=\sum_{n\geq0}(Da_n)Z^n.
\]
This is a formal construction: no evaluation at a scalar is being made.
It is a derivation since each coefficient of a product is a finite sum.
The formal derivations $\DD$ and $\partial_Z$ commute, and their sum
$\DD+u\partial_Z$ is a derivation whenever $u$ is a coefficient scalar.
Commutation with that sum need not hold when $Du\ne0$; the elementary
commutator records the changing scale.

\begin{proposition}[Two-derivative chain rule]\label{prop:chain}
For $P(Z)\in\K[Z]$ and $h\in\K$,
\begin{equation}\label{eq:two-chain}
 D(P(h))=(\DD P)(h)+(\partial_ZP)(h)Dh.
\end{equation}
For $F(Z)=\sum_{n\geq0}a_nZ^n$, suppose the evaluated family
$(a_nh^n)_{n\geq0}$ is summable. Then the two families defining
$(\DD F)(h)$ and $(\partial_ZF)(h)$ are summable as well, and
\eqref{eq:two-chain} holds with $F$ in place of $P$.
\end{proposition}
\begin{proof}
For a polynomial, differentiate each term using Leibniz' rule:
\[
 D(a_nh^n)=(Da_n)h^n+n a_nh^{n-1}Dh.
\]
There are finitely many terms, so summing proves \eqref{eq:two-chain}.

For a series, first consider $h\ne0$. Multiplying each member of a
summable family by its ordinary integer index preserves summability:
it does not enlarge any support. Multiplication of the whole family by
the fixed scalar $h^{-1}$ also preserves Hahn summability. Thus
\[
 (n a_nh^{n-1})_{n\geq1}
 =(n h^{-1}(a_nh^n))_{n\geq1}
\]
is summable. Its scalar multiple by $Dh$ is summable too. Strong
additivity says that $(D(a_nh^n))_n$ is summable. Subtracting the
already controlled variable-derivative family proves summability of
$((Da_n)h^n)_n$. The termwise product rule can now be summed, giving
\eqref{eq:two-chain}. When $h=0$, the evaluated series and both formal
derivative evaluations have at most one nonzero term each, and the
identity reduces to $Da_0=Da_0+a_1D0$.
\end{proof}

The support argument is part of the theorem, not an interchange of
analytic limits. An arbitrary sum of two families can be summable through
cancellation even when its separate families are not. Here the special
power-series form controls the variable-derivative family first; only
then is the split justified. This gives a stronger statement than merely
assuming three independent summability certificates, but still does not
assert that an arbitrary formal series is evaluable at an arbitrary scalar.

Here is a finite check exposing both terms. Let
$P(Z)=\omega Z+tZ^2$ and $h=\omega^2$. Direct substitution gives
$P(h)=2\omega^3$, so $D(P(h))=6\omega^2$. On the right,
\[
 (\DD P)(h)=h-t^2h^2=0,
 \qquad
 (\partial_ZP)(h)Dh=(\omega+2th)(2\omega)=6\omega^2.
\]
The cancellation in the coefficient term is part of this example, not a
reason to omit coefficient differentiation in general.

\subsection{A fine derivative does not determine an intrinsic one}

The additive projections in \eqref{eq:three-parts} give especially clear
counterexamples. For every $x\in\No$ and every finite increment $h$,
\[
 \fin(x+h)=\fin(x)+h,\qquad
 \pipart(x+h)=\pipart(x).
\]
Therefore the fine derivative of $x\mapsto\fin(x)$ is $1$ everywhere,
and that of $x\mapsto\pipart(x)$ is $0$ everywhere. Indeed these difference
quotients are exact for all nonzero $\abs h<1$, so there is no subtle
limit calculation.

At $x=\omega$, however,
\[
 D(\fin(\omega))=D0=0,\qquad
 D(\pipart(\omega))=D\omega=1.
\]
Thus a blanket identity $D(F(x))=F'_{\mathrm{fine}}(x)Dx$ is false for
arbitrary fine-differentiable functions. The missing hypothesis is a
compatibility structure for $F$; fine differentiability alone is not one.

These examples also explain why a fine derivative equal to zero need not
force a global function to be constant. The purely infinite projection is
locally constant for finite perturbations but is globally nonconstant.
This is a function-theoretic issue, whereas $Df=0$ for a scalar means
exactly $f\in\C$ in the intrinsic theory. Both statements are true, and
conflating them creates a false contradiction.

\subsection{Ordinary analytic parameters}

One can safely use coefficientwise differentiation and coefficientwise
contour integrals in an ordinary parameter without choosing any intrinsic
scalar derivation. After a derivation is chosen, exchanging it with such
an operation needs its own support and domain theorem. The finite example
\eqref{eq:deriv-separator} already shows why no unqualified exchange rule
should be assumed.

For a fixed ordinary point $z\in U$, strong additivity does compute
$D(F(z))=\sum_\gamma f_\gamma(z)D(t^\gamma)$. This is a pointwise scalar
identity. Turning it into a common-domain coherent family, differentiating
again in $z$, or integrating it uniformly along a path requires control of
the coefficient-family supports and of the domain. The distinct coefficient
rings documented in the repository remain distinct after $D$ is added.

\section{A genuine bridge: omega-series and composition}
\label{sec:composition}

There is an established setting in which intrinsic differentiation really
can be recovered as differentiation of a represented function. It is not
all of $\No$ with an arbitrary choice of functions.

\subsection{The restricted composition theorem}

Let $\mathscr W=\R\langle\!\langle\omega\rangle\!\rangle$ denote the
omega-series: the smallest subfield of $\No$ containing $\R$ and $\omega$
and closed under exponential, positive logarithm, and set-indexed summable
families. It is a proper class, not our set-sized workspace $\HH$.
Berarducci--Mantova construct a composition $f\circ x$ for
$f\in\mathscr W$ and positive infinite $x$. Their results identify
$(Df)\circ x$ with the fine derivative of $x\mapsto f\circ x$, and give
the compatible chain rule when composition stays in the omega-series
setting. Their Taylor theorem states that, for sufficiently small $h$,
\begin{equation}\label{eq:omega-Taylor}
 f\circ(x+h)=\sum_{n\geq0}\frac{(D^nf)\circ x}{n!}h^n.
\end{equation}
The sum is support-controlled, and the permitted size of $h$ depends on
$f$ and $x$ \cite[Corollaries 7.6--7.7, Theorems 7.14 and 8.2]{BMGerms}.

These are imported composition and Taylor results. They do not follow
merely from the derivation axioms in Theorem~\ref{thm:BM}. In particular,
$f\circ x$ is an additional operation, not the result of substituting into
an arbitrary normal form without a summability proof.

\subsection{How the two derivatives agree}

For $f=\log\omega$, the represented function is $x\mapsto\log x$ and
$Df=1/\omega$. Composition sends the latter to $1/x$, exactly its fine
function derivative. For $f=\omega^3$, the two descriptions similarly give
$3x^2$. Equation~\eqref{eq:omega-Taylor} explains this agreement in a
much larger class than finite elementary expressions.

This is a different construction from the formal polynomial
$P(Z)=\omega Z$. In that formal polynomial, $\omega$ is a coefficient
held fixed by $\partial_Z$. In the omega-series $f=\omega^2$, both
occurrences belong to the distinguished compositional scale, and
$f\circ x=x^2$. A representation must record which interpretation is
intended before any chain rule can be chosen.

The local equality of derivative operations also does not say that every
surreal scalar is an omega-series or that every primitive belongs to that
class. The full-field surjectivity theorem and the restricted composition
theorem have different domains and must be applied separately.

\subsection{A boundary, not a global substitution engine}

Berarducci--Mantova also prove that their simplest derivation on all of
$\No$ cannot be compatible with a global composition satisfying the
composition and compatibility axioms of that paper
\cite[Theorem 8.4]{BMGerms}. This is a specific incompatibility result, not
a prohibition on every conceivable partial substitution operation.

The practical lesson is to retain at least three different input kinds:
a formal expression with named coefficient constants, an omega-series with
a distinguished compositional scale, and a general scalar with an
intrinsic derivation. They can overlap as values while supporting different
operations. Identifying their carriers in memory does not justify identifying
their semantic interfaces.

\section{An oscillatory differential extension outside the surcomplex field}
\label{sec:PV}

The equation $Dy=iy$ is consistent as an ordinary differential-algebraic
equation. Its lack of nonzero surcomplex solutions means that a differential
extension providing a solution cannot in general be represented by adding
another ordinary surcomplex scalar.

\subsection{A set-sized construction}

Let $F=\C(X)$ with $DX=1$, and adjoin a transcendental indeterminate $T$.
On $E=F(T)$ define
\begin{equation}\label{eq:PV-deriv}
 DX=1,\qquad DT=iT,
\end{equation}
then extend by the polynomial and quotient rules. These prescriptions define
a derivation because $X$ and $T$ are algebraically independent. Thus $T$ is
a nonzero solution of the required equation in a perfectly valid set-sized
differential field.

The usual Picard--Vessiot perspective asks for an extension generated by a
fundamental solution without adjoining new differential constants
\cite{Singer}. We verify that property here rather than assume it.

\begin{lemma}\label{lem:rational-logder}
For $r(X)\in\C(X)^\times$, the identity $r'/r=c$ with a nonzero constant
$c\in\C$ is impossible.
\end{lemma}
\begin{proof}
Write $r=p/q$ with nonzero complex polynomials. If $d=\deg p$ and
$e=\deg q$, expansion at infinity gives
\[
 \frac{r'}r=\frac{p'}p-\frac{q'}q
 =\frac{d-e}{X}+O(X^{-2}).
\]
This is an identity of formal Laurent expansions in $1/X$, whose constant
term is zero, so it cannot equal a nonzero constant. Equivalently, one can
compare degrees after clearing denominators in
$p'q-pq'=cpq$.
\end{proof}

\begin{proposition}\label{prop:PVconstants}
The differential constant field of $E=\C(X)(T)$ with
\eqref{eq:PV-deriv} is exactly $\C$.
\end{proposition}
\begin{proof}
Expand a rational function $R(X,T)$ at $T=0$:
\[
 R(X,T)=\sum_{n\geq n_0}r_n(X)T^n\in\C(X)((T)).
\]
This is an injective formal Laurent expansion; it asserts no analytic
convergence. The derivation on the Laurent-series field defined by
\[
 D\left(\sum_n r_n(X)T^n\right)
   =\sum_n\bigl(r_n'(X)+in r_n(X)\bigr)T^n
\]
is compatible with the rational subfield and with
\eqref{eq:PV-deriv}. If $DR=0$, each coefficient satisfies
$r_n'+inr_n=0$. For $n\ne0$, Lemma~\ref{lem:rational-logder} forces
$r_n=0$. For $n=0$, the equation says $r_0'=0$, hence $r_0\in\C$.
Injectivity of the expansion now gives $R\in\C$.
\end{proof}

The field $E$ is generated by $T$, a one-dimensional fundamental solution,
and has no new constants, so it is a rank-one Picard--Vessiot field for
$DY=iY$ over $F$. Its differential automorphisms over $F$ are exactly
\[
 T\longmapsto cT,\qquad c\in\C^\times.
\]
Indeed every such scaling is an automorphism commuting with $D$; conversely,
the ratio of two nonzero solutions is constant. Its differential Galois
group is therefore the multiplicative group $\mathbb G_m(\C)$.

\subsection{Why no differential embedding is possible}

The base field $F$ embeds in $\K$ by $X\mapsto\omega$ and by fixing
$\C$. This is a differential embedding, since $D\omega=1$. It cannot
extend to a differential embedding of $E$: the image of $T$ would be
nonzero and satisfy $DY=iY$, contradicting
Proposition~\ref{prop:no-oscillation}.

Thus a field can be algebraically closed, contain the required base
coefficients, and possess additive primitives, yet still fail to contain
this simplest oscillatory differential extension. The role of $T$ is not
that of an algebraic root missing from $\No[i]$; it is a new differential
transcendental. No change of normal-form encoding inside the same field
can make it exist.

A symbolic constructor for such an extension is mathematically reasonable,
but it must announce a change of ambient differential field. Calling its
new element a surcomplex number would erase the distinction proved here.
We use only the set-sized base $\C(X)$ in this construction; no blanket
application of Picard--Vessiot existence theory to a proper-class base is
needed.

\section{Computational and formalization consequences}
\label{sec:implementation}

\subsection{An executable exact core}

For finite rational Hahn polynomials with Gaussian-rational coefficients,
represent an input by a finite map
\[
 q\in\Q\ \longmapsto\ c_q\in\Q(i),
\]
with zero coefficients deleted. Exact addition and multiplication use the
usual sparse coefficient rules. Proposition~\ref{prop:Hahn-deriv} gives a
terminating derivative operation. For a finite input $f$, additive
integration inside $\HH$ succeeds exactly when the exponent-$1$
coefficient is zero; otherwise the formula
\eqref{eq:Hahn-log-integral} gives an ambient primitive using one explicit
$\log\omega$ constructor.

For $Dy=qy$, two tests must be offered separately. The \emph{workspace}
test rejects every nonzero term below exponent $1$, and rejects a
nonreal exponent-$1$ coefficient. With Gaussian-rational coefficients, a
real exponent-$1$ coefficient is automatically rational. The remaining
positive-valuation primitive $k$ certifies the answer
$t^r\EE(k)$. The \emph{ambient} test examines the imaginary-part polynomial
$b$: it succeeds precisely when $b=0$ or every exponent in $\supp b$ is
strictly greater than $1$.

The output $t^r\EE(k)$ is a finite exact expression denoting an infinite
Hahn series, not a finite list of all its coefficients. Its certificate
consists of the rational exponent $r$, the finite polynomial $k$, its
positive support, and the exact equality $-rt+Dk=q$.
For an ambient-only result, a displayed primitive and real exponential may
supply a finite expression even though the chosen backend cannot expand it
in its workspace.

\subsection{Return a reason, not just a failure flag}

\begin{center}
\small
\begin{tabular}{@{}p{0.30\textwidth}p{0.60\textwidth}@{}}
\toprule
Result & Mathematical meaning\\
\midrule
Workspace solution & A certified nonzero solution lies in the declared Hahn field.\\[3pt]
Ambient-only solution & A solution exists in $\No[i]$, but the workspace obstruction is nonzero.\\[3pt]
Ambient phase obstruction & The purely infinite part of a phase primitive is nonzero; no nonzero surcomplex solution exists.\\[3pt]
Undetermined representation & The supplied description does not support a certified support or primitive calculation. This is not a nonexistence theorem.\\
\bottomrule
\end{tabular}
\end{center}

The first three outcomes are decidable for the finite exact core just
described. For arbitrary computably described Hahn supports, they are not
advertised here as total decision procedures. The semantic formula
$\Omega(q)=\pipart(J\Imm q)$ remains exact without promising an algorithm
for $J$ or for equality of arbitrary represented scalars. This matches the
capability-tier discipline already used by the repository \cite{RepoCAS}.

\subsection{The delivered verification program}

The archive includes a Python program using exact rational arithmetic and
SymPy for finite symbolic identities. It checks rational monomial
primitives, product and quotient rules, the two-derivative polynomial
identity, several phase certificates, the corrected normalized integration
identity, and finite polynomial inverses of $D-i$. It also implements the
finite-input workspace and ambient existence tests above and checks their
classification on representative inputs.

These checks do not enumerate arbitrary Hahn supports or implement $D$ on
all surreal exponents. A finite Taylor check is explicitly only a formal
coefficient check up to its stated order. The general support arguments
remain the mathematical proofs in this article. The verification output
records the actual test count and software versions, separately from the
LaTeX build.

\subsection{A formalization boundary}

A clean proof-assistant development can start with a set-sized ordered
field $F$, a derivation, a chosen infinite scale $w$ with $Dw=1$, real
constants, and specified local exponential/logarithm operations with their
support theorems. The bounded-derivative argument, complexification,
phase decomposition, and first-order criterion then become lemmas over
explicit hypotheses. Surjectivity of the derivation is a genuine additional
hypothesis, not something inferred from real closedness.

For the explicit rational Hahn workspace, the structural theorem to prove
first is \eqref{eq:Hahn-D}. The primitive-image and logarithmic-image
proofs then reduce to support translation, leading-term factorization,
and the infinitesimal exponential theorem. These provide a tractable
independent target before a full formal construction of the
Berarducci--Mantova derivation is available.

The interfaces should keep $D$, $\partial_Z$, $\DD$, coefficientwise
analytic differentiation, and fine differentiability separate. Evaluation
is a map with a support certificate, and Proposition~\ref{prop:chain} is a
compatibility theorem for that map. It is not a coercion identifying the
operators. Likewise, an exponential operation should carry the domain
predicate $\Imm z\in\OO$ when it promises the intrinsic differential law.

No Lean source or machine-checked formalization is included in this
package. The preceding paragraph specifies proof obligations and a possible
order of development; it does not claim those modules already exist in the
repository or in a library.

\section{Conclusions and further questions}
\label{sec:conclusion}

The intrinsic differential field and the analytic theories in the
repository are compatible only after their distinct operations and domains
are made explicit. The decisive structural facts can be summarized without
using a global complex exponential:
\[
 \ker D=\C,\qquad D\K=\K,\qquad
 \Dlog(\K^\times)=\No+iD(\mm).
\]
The last image is proper. Its obstruction is not lack of algebraic closure
or of additive integration, but the purely infinite part of an imaginary
primitive. The same principle identifies the largest possible argument
domain for an intrinsically exponential-compatible value and removes all
nonreal characteristic roots from the solution space of a
constant-coefficient equation.

At the workspace level, two even sharper identities are available:
\[
 D(\HH)=\{f\in\HH:[t]f=0\},\qquad
 \Dlog(\HH^\times)=\Q t+\{g:g=0\text{ or }v(g)>1\}.
\]
They give exact, auditable tests rather than heuristic simplification rules.
They also show why derivative closure, primitive closure, and exponential
closure must be separate capabilities.

Several natural continuations remain beyond the results proved here.
One is to classify logarithmic derivatives in multi-scale Hahn workspaces
stable under $D$, where support shifts no longer have the uniform rational
form of \eqref{eq:Hahn-D}. Another is to determine useful classes of
variable-coefficient matrix equations whose fundamental matrices remain in
$\No[i]$, distinguishing genuine phase obstructions from workspace
obstructions. A third is to build explicitly typed analytic/coefficient
interfaces that prove exchange rules with contour functionals under their
correct common-domain and summability hypotheses. These are directions for
further development, not assertions that the corresponding questions are
unresolved in the literature.

For the present documentation gap, the essential bridge is now concrete:
\emph{an intrinsic derivative of a scalar, an external derivative of a
function, and the substitution that might connect them are three pieces of
structure, not one.}

\appendix
\section{Notation and theorem dependencies}
\label{app:notation}

\begin{longtable}{@{}p{0.23\textwidth}p{0.69\textwidth}@{}}
\toprule
Notation & Meaning\\
\midrule
\endhead
$\No$, $\K$ & Surreal real field and its complexification $\No[i]$; class-sized carriers.\\[3pt]
$D$ & Fixed normalized Berarducci--Mantova derivation; $D\omega=1$.\\[3pt]
$\OO$, $\mm$ & Finite real surreals and real infinitesimals.\\[3pt]
$\PP$ & Purely infinite real normal forms, including zero.\\[3pt]
$\st$, $\ct$ & Standard part on finite elements; coefficient of $t^0$ on arbitrary normal forms. They agree only on the finite domain.\\[3pt]
$\pipart$, $\fin$ & Purely infinite and finite-part projections in the additive normal-form decomposition.\\[3pt]
$t^\gamma$ & Conway monomial $\omega^{-\gamma}$, not a generic surreal exponential power.\\[3pt]
$\EE$, $\LL$ & Support-defined exponential on $\mm_\K$ and logarithm on $1+\mm_\K$.\\[3pt]
$\Dlog$ & Logarithmic derivative $y\mapsto Dy/y$.\\[3pt]
$\Omega$ & Purely infinite part of a real primitive of the imaginary coefficient.\\[3pt]
$J$ & Unique primitive with zero ordinary constant coefficient.\\[3pt]
$\PS$ & Additive exponential domain $\No+i\OO$, not a subring.\\[3pt]
$\HH$ & Set-sized rational Hahn field $\C((t^\Q))$.\\[3pt]
$K_\Gamma$ & Full Hahn workspace with a set-sized divisible exponent group.\\[3pt]
$\partial_Z$, $\DD$ & Formal variable differentiation and intrinsic coefficient differentiation on $\K[[Z]]$.\\[3pt]
$\mathscr W$ & Proper-class omega-series field with the specified partial composition theory.\\
\bottomrule
\end{longtable}

The logical dependency chain of the central criterion is short:
\[
 \begin{gathered}
 \text{BM structure}\\
 \Downarrow\\
 \text{complexification and local exponential calculus}\\
 \Downarrow\\
 \text{unit-phase decomposition}\ \Longrightarrow\
 \text{first-order criterion}.
 \end{gathered}
\]
The bounded-derivative lemma gives an independent quick proof of
$Dy=iy\Rightarrow y=0$. The exact sequence and maximal exponential domain
then follow from the full criterion. The rational Hahn calculations use
only the rational monomial restriction of $D$ together with the local
exponential. The composition theorem and the Picard--Vessiot background
are not dependencies of the phase criterion.

\section{Verification and reproduction boundary}
\label{app:verification}

The article is a standalone LaTeX source with an internal bibliography.
Build it with a standard TeX installation using
\begin{verbatim}
latexmk -pdf -interaction=nonstopmode article.tex
\end{verbatim}
No external figure, bibliography database, shell escape, or network access
is needed to compile it. The program is separate:
\begin{verbatim}
python verify_examples.py
\end{verbatim}
The accompanying \code{verification.txt} records its executed result.
The finite-input classifications are mathematical consequences of
Theorem~\ref{thm:Hahn-logder} and Proposition~\ref{prop:Hahn-ambient}; the test cases
illustrate those consequences and guard against sign or scale errors.

The repository audit was read-only and pinned to the revision recorded in
Section~\ref{sec:gap}. No repository source was edited, no package was
claimed to be executed, and no source manuscript was silently merged into
this article. The audit log in the archive identifies the inspected paths
and the exact scope of the gap assessment.

\clearpage
\begin{thebibliography}{99}\raggedright
\addcontentsline{toc}{section}{References}
\small

\bibitem{BM}
A. Berarducci and V. Mantova,
\emph{Surreal numbers, derivations and transseries},
Journal of the European Mathematical Society \textbf{20} (2018), 339--390.
\href{https://doi.org/10.4171/JEMS/769}{doi:10.4171/JEMS/769}.
Preprint: \arxiv{1503.00315v3}.
The structural inputs used here are Theorems 6.30 and 7.7.

\bibitem{BMGerms}
A. Berarducci and V. Mantova,
\emph{Transseries as germs of surreal functions},
Transactions of the American Mathematical Society \textbf{371} (2019),
3549--3592.
\href{https://doi.org/10.1090/tran/7428}{doi:10.1090/tran/7428}.
Preprint: \arxiv{1703.01995v3}.
The restricted composition, Taylor, and incompatibility results are used
only in Section~\ref{sec:composition}.

\bibitem{Gonshor}
H. Gonshor,
\emph{An Introduction to the Theory of Surreal Numbers},
London Mathematical Society Lecture Note Series, vol.~110,
Cambridge University Press, 1986.
Publisher catalogue:
\url{https://www.cambridge.org/core/books/an-introduction-to-the-theory-of-surreal-numbers/312AE504A3E88E804054BFB390446374}.

\bibitem{Neumann}
B. H. Neumann,
\emph{On ordered division rings},
Transactions of the American Mathematical Society \textbf{66} (1949),
202--252.
\href{https://doi.org/10.1090/S0002-9947-1949-0032593-5}{doi:10.1090/S0002-9947-1949-0032593-5}.
The positive-support summability lemma is used in its standard
additive-exponent form.

\bibitem{BG}
O. Bournez and Q. Guilmant,
\emph{Surreal fields stable under exponential, logarithmic, derivative and
anti-derivative functions}, preprint (2022), \arxiv{2211.08396}.
Cited for further structural workspace results, not as a dependency of
Theorem~\ref{thm:localization}.

\bibitem{Singer}
M. F. Singer,
\emph{Introduction to the Galois theory of linear differential equations},
in M. A. H. MacCallum and A. V. Mikhailov (eds.),
\emph{Algebraic Theory of Differential Equations},
London Mathematical Society Lecture Note Series, vol.~357,
Cambridge University Press, 2009, pp.~1--82.
Preprint: \arxiv{0712.4124v2}.

\bibitem{RepoMap}
\repo\ contributors,
\emph{The surreal and surcomplex reports}, documentation catalogue,
\code{docs/README.md}, inspected 21 September 2026 at revision
\code{e260237db9b71da8b74a0c13c8e6355119091100}.
\url{https://github.com/VladimirReshetnikov/Surreal/blob/e260237db9b71da8b74a0c13c8e6355119091100/docs/README.md}.

\bibitem{RepoAnalysis}
\repo\ contributors,
\emph{Surcomplex analysis: holomorphic functions on K = No[i]},
\code{docs/surcomplex/analysis/README.md} and the opening foundations
of \code{article.tex}; same inspected revision as \cite{RepoMap}.
\url{https://github.com/VladimirReshetnikov/Surreal/tree/e260237db9b71da8b74a0c13c8e6355119091100/docs/surcomplex/analysis}.

\bibitem{RepoTrig}
\repo\ contributors,
\emph{Trigonometry on the Surcomplex Plane},
\code{docs/surcomplex/trigonometry/README.md}; same inspected revision
as \cite{RepoMap}. In particular, its stated boundary excludes choosing an
intrinsic derivation or classifying differential-equation solutions.
\url{https://github.com/VladimirReshetnikov/Surreal/blob/e260237db9b71da8b74a0c13c8e6355119091100/docs/surcomplex/trigonometry/README.md}.

\bibitem{RepoFound}
\repo\ contributors,
\emph{Foundations for surreal and surcomplex mathematics},
\code{docs/foundations-and-computation/foundations/README.md};
same inspected revision as \cite{RepoMap}.
\url{https://github.com/VladimirReshetnikov/Surreal/blob/e260237db9b71da8b74a0c13c8e6355119091100/docs/foundations-and-computation/foundations/README.md}.

\bibitem{RepoCAS}
\repo\ contributors,
\emph{Surreal and Surcomplex Numbers in Symbolic Computer Algebra},
\code{docs/foundations-and-computation/computer-algebra/README.md};
same inspected revision as \cite{RepoMap}. The consulted excerpt covered
its capability, scope and prototype descriptions, not a complete execution
audit of its source packages.
\url{https://github.com/VladimirReshetnikov/Surreal/blob/e260237db9b71da8b74a0c13c8e6355119091100/docs/foundations-and-computation/computer-algebra/README.md}.

\end{thebibliography}
\end{document}
